% riess_schmidt_perlmutter_ML_short.tex
% Short newsletter version of riess_schmidt_perlmutter_ML.tex
% Thomas J. Sargent, 2026

\documentclass[11pt]{article}

\usepackage{amsmath,amssymb}
\usepackage{booktabs}
\usepackage[margin=1.15in]{geometry}
\usepackage{microtype}
\usepackage{hyperref}
\usepackage{natbib}
\usepackage{parskip}

\hypersetup{colorlinks=true, linkcolor=blue, citecolor=blue, urlcolor=blue}

\title{Fifty Supernovae and One Combination of Two Parameters\\[2pt]
       \large The Accelerating Universe as Statistical Inference, 1988--1999}
\author{Thomas J.\ Sargent\thanks{%
  In 2011 my good friend Christopher A.\ Sims and I shared the Nobel Prize in
  economics.  That good luck let us meet and become friends with winners of
  prizes in physics and other fields.  Adam Riess had majored partly in
  economics as an undergraduate.  He and Brian Schmidt have since taught me
  the nuts and bolts of the computer programs they used to climb the
  likelihood function that led to their discovery.
  This paper is one of five that read episodes in the
  history of physics and astronomy as exercises in statistical inference and
  machine learning.  \citet{sargent2025sources} sets out the theme: what is
  now called artificial intelligence names practices that scientists have used
  for centuries.  The other four treat Kepler and Newton, Einstein's photoelectric equation,
  Pauli's exclusion principle, and the exoplanet discovery of Mayor and
  Queloz.}}
\date{April 2026}

\begin{document}
\maketitle

\begin{abstract}
\noindent
In 1998 two teams reported that distant supernovae are fainter than a
decelerating universe allows, and concluded that cosmic expansion is speeding
up.  The inference is a weighted nonlinear regression with two parameters of
interest and a nuisance intercept.  I run the Fisher-information calculation
that prices it.  The result reproduces the published error bar from the
observing design alone, and shows that the supernovae measured one combination
of the two density parameters and left the orthogonal combination nearly free.
The famous numbers assume a flat universe.  Ruling out a decelerating universe
needs a prior the supernovae do not supply.  I also set out the episode's
genuine out-of-sample test, which came six years later and had no free
parameters, and I ask what the whole affair says about theory-free discovery.
The derivations and the full data are in a longer companion version.
\end{abstract}

%======================================================================
\section{What Was Measured}

In 1998 two teams looked at exploding stars billions of light-years away and
found them fainter than they should have been.  Fainter means further.
Further, at a given redshift, means the universe has expanded more since the
light left than anyone had allowed for.  The expansion is speeding up, and
something is pushing it.  Nobody knows what.

The inference from those photometric measurements to that conclusion is a
regression, and this note is about its statistical anatomy: what was
identified, what was assumed, and how much of the famous result the data
actually carry.

Start with the measurement.  A \emph{standard candle} is an object whose
intrinsic brightness you know.  Knowing it, you convert observed brightness
into distance, since flux falls as the inverse square.  Collect distances and
redshifts for objects near and far and you have traced the expansion history
of the universe.

Type Ia supernovae are thermonuclear explosions of white dwarfs in binary
systems.  Their peak luminosities are nearly, but not quite, uniform.  The
account current in the 1990s had the white dwarf accreting until it neared the
Chandrasekhar limit, so that a near-universal exploding mass gave a
near-universal brightness.  That account no longer commands consensus.
Sub-Chandrasekhar detonations and white-dwarf mergers are now taken seriously,
and the progenitor population may be mixed.  The uniformity is an empirical
fact, not a derived one, and it was empirical in 1998 too.

Magnitudes are a logarithmic brightness scale on which five magnitudes is a
factor of one hundred in flux, and larger means fainter.  Raw Type~Ia peak
magnitudes scatter by about $0.4$ magnitudes, a $20\%$ distance error.  The
cosmological signal these teams were chasing is about $0.25$ magnitudes.
Before 1993 the noise was larger than the thing being measured.

Two groups went after it anyway.  The Supernova Cosmology Project, led by Saul
Perlmutter at Berkeley, began systematic high-redshift searches in the late
1980s and reported 42 supernovae at redshift $z = 0.18$ to $0.83$
\citep{Perlmutter1999}.  The High-Z Supernova Search Team, assembled around
Brian Schmidt with Adam Riess writing the analysis pipeline, reported 16
high-redshift events at $z = 0.16$ to $0.97$ together with 34 nearby ones
\citep{Riess1998}.  Both concluded that the expansion is accelerating.
Perlmutter took half the 2011 Nobel Prize; Riess and Schmidt shared the other
half.

%======================================================================
\section{The Regression}

General relativity, plus the assumption that the universe looks the same in
every direction and from every place, gives the expansion rate as
\begin{equation}\label{eq:friedmann}
  H^2(z) = H_0^2\bigl[\Omega_M(1+z)^3 + \Omega_k(1+z)^2 + \Omega_\Lambda\bigr],
\end{equation}
where $\Omega_M$ is the density of matter, $\Omega_\Lambda$ the density of the
cosmological constant, and $\Omega_k$ spatial curvature, all in units that
make them sum to one.  A \emph{flat} universe has $\Omega_k = 0$.  Integrating
\eqref{eq:friedmann} along the light path gives the luminosity distance
$d_L(z; \Omega_M, \Omega_\Lambda)$, and taking logarithms gives the
\emph{distance modulus} $\mu = 5\log_{10}(d_L/\text{Mpc}) + 25$, which is what
a supernova's apparent magnitude measures once its intrinsic magnitude is
known.

The problem is well behaved for two reasons.  The functional form is fixed,
exactly, by physics established seventy years earlier.  And the
Hubble constant $H_0$ and the intrinsic magnitude $M_B$ enter only through a
single combination,
\begin{equation}\label{eq:calM}
  \mathcal{M} = M_B + 5\log_{10}\!\left(\frac{c}{H_0\,\text{Mpc}}\right) + 25,
\end{equation}
so the measurement equation is
\begin{equation}\label{eq:meas}
  m_{B,i} = \mathcal{M}
    + 5\log_{10}\tilde d_L(z_i;\Omega_M,\Omega_\Lambda) + \varepsilon_i,
  \qquad \varepsilon_i \sim \mathcal{N}(0,\sigma_i^2),
\end{equation}
where $\tilde d_L$ is the dimensionless part of the distance that depends only
on the density parameters.

$\mathcal{M}$ is a textbook nuisance parameter.  It shifts the whole Hubble
diagram vertically and leaves its shape alone.  Neither $H_0$ nor $M_B$ can be
recovered from the diagram, because any change in one is offset exactly by a
change in the other.  Neither needs to be.  Since $\mathcal{M}$ enters
linearly it can be concentrated out analytically at each candidate
$(\Omega_M, \Omega_\Lambda)$, leaving a profile criterion in two parameters.

Millikan faced the same structure in 1916 measuring Planck's constant.
The contact voltage between his photoemitting metal and his collector shifted
every reading by the same amount and left the slope he wanted untouched.  The
parallel is exact as far as it goes.  Section~\ref{sec:fisher} shows where it
stops.

Under Gaussian errors, maximising the likelihood means minimising
\begin{equation}\label{eq:chi2}
  \chi^2 = \sum_{i}
  \frac{\bigl[m_{B,i} - \mathcal{M}
    - 5\log_{10}\tilde d_L(z_i;\Omega_M,\Omega_\Lambda)\bigr]^2}{\sigma_i^2},
\end{equation}
which is weighted least squares.  Unlike Millikan's line, $\tilde d_L$ has no
closed form, so both groups minimised \eqref{eq:chi2} on a grid, computing the
integral numerically at each point.

\subsection*{The easy case, in the preceding note}

The preceding note works through Millikan's 1916 measurement of Planck's
constant, which has this anatomy and one simplification.  A contact voltage
between his photoemitting metal and his collector shifted every reading by the
same amount, as $\mathcal{M}$ shifts every point of the Hubble diagram, and a
slope carried the physics that a level could not touch.

Millikan wanted one number.  Removing his nuisance parameter left that number
identified, and five observations measured it to under one percent.  The
Friedmann model wants two, and removing $\mathcal{M}$ leaves them entangled
with each other, so a second degeneracy survives the first.
Section~\ref{sec:fisher} is about that second degeneracy.  It belongs to the
geometry of the problem rather than to any instrument, and no supernova
measurement removes it.

%======================================================================
\section{The Feature That Made It Possible}

None of this helps while the noise is $0.4$ magnitudes and the signal is
$0.25$.

\citet{Phillips1993} found the fix.  Plotting the rate at which a supernova
fades after its peak against how bright that peak was, he found a tight
relation: slower-declining supernovae are brighter.  Writing $\Delta m_{15}$
for the drop in $B$-band magnitude over the fifteen days after peak,
\begin{equation}\label{eq:phillips}
  M_B^{\rm peak} \approx M_B^{\rm ref} + \alpha\,(\Delta m_{15} - 1.1),
  \qquad \alpha \approx 0.75~\text{mag},
\end{equation}
with $\alpha$ here the \citet{Hamuy1996} calibration.  Correcting for it cuts
the dispersion from $0.4$ to about $0.15$ magnitudes.

That factor of $2.7$ in the noise is a factor of seven in the number of
supernovae needed for a given precision.  Against the open matter-only model
the signal at $z \approx 0.5$ is $0.22$ magnitudes, and detecting it at three
standard errors takes about thirty supernovae at $\sigma = 0.4$ and about four
at $\sigma = 0.15$.  Both teams had roughly the number of events that the
pre-Phillips dispersion would have required just to see the effect, with none
left over to separate it from anything else.

\subsection*{The half of the discovery that came from data}

The story so far credits theory.  General relativity fixed the functional
form, three parameters remained, and fifty supernovae sufficed where a
theory-free learner would have needed thousands.  Half the discovery does not
fit that account.

$\Delta m_{15}$ came out of data.  No theory of Type Ia explosions predicted
that decline rate tracks peak luminosity, and Phillips did not derive it.  He
plotted decline rate against absolute magnitude for nine nearby supernovae and
fitted a line.  The relation has since been partly rationalised---more
radioactive nickel yields both a brighter peak and slower-cooling ejecta---and
it has never been deduced.  Thirty years on it remains an empirical
regularity, recalibrated with every new survey.  A theorist in 1993 holding
the correct progenitor model could not have written \eqref{eq:phillips}, and a
theorist today still cannot.

So the supernova discovery was a hybrid.  A deductively specified hypothesis
class was bolted to an inductively discovered feature, and neither half would
have done.  Without Friedmann there is no curve to fit and no meaning for the
fitted numbers.  Without Phillips the scatter is twice the signal and no
amount of theory recovers it.  The one step where a data-driven procedure did
the indispensable work is the one step where no theory was available to do it
instead.

The pattern has continued.  Modern supernova cosmology runs on SALT2 and
SALT3, light-curve models trained on data rather than derived from explosion
physics, with training sets that grow every survey.  Feature extraction in
this field has become more machine-learned over time.  The hypothesis class
has stayed Friedmann's.

%======================================================================
\section{The Signal}

Table~\ref{tab:mu} gives the distance modulus predicted by three cosmologies.

\begin{table}[ht]
\centering
\caption{Predicted distance modulus for
$H_0 = 70$~km\,s$^{-1}$\,Mpc$^{-1}$, by numerical integration.  The last two
columns give the excess faintness of flat $\Lambda$CDM over the open
matter-only model and over Einstein--de Sitter.}
\label{tab:mu}
\medskip
\begin{tabular}{lcccrr}
\toprule
$z$ & Flat $\Lambda$CDM & Einstein--de Sitter & Open
  & $\Lambda-$open & $\Lambda-$EdS \\
 & $(0.3,\,0.7)$ & $(1.0,\,0.0)$ & $(0.3,\,0.0)$ & & \\
\midrule
0.10 & 38.32 & 38.21 & 38.25 & $+0.07$ & $+0.11$ \\
0.25 & 40.50 & 40.27 & 40.36 & $+0.14$ & $+0.24$ \\
0.50 & 42.26 & 41.86 & 42.05 & $+0.22$ & $+0.40$ \\
0.75 & 43.33 & 42.82 & 43.08 & $+0.25$ & $+0.51$ \\
1.00 & 44.10 & 43.50 & 43.84 & $+0.26$ & $+0.60$ \\
\bottomrule
\end{tabular}
\end{table}

At $z \approx 0.5$ the models differ by $0.22$ to $0.40$ magnitudes, against a
per-supernova error of about $0.22$ magnitudes.  Each event carries roughly
one standard deviation of signal.  Both teams observed a weighted mean excess
of about $+0.25 \pm 0.06$ magnitudes.

Naive stacking would put forty such events at $\sqrt{40} \approx 6$ standard
errors.  Both teams claimed about three.  The gap is the subject of the next
section, and it is the most interesting statistical feature of the episode.

The published estimates, all of which impose flatness, are in
Table~\ref{tab:results}.

\begin{table}[ht]
\centering
\caption{Published estimates, all imposing $\Omega_M + \Omega_\Lambda = 1$.
Statistical uncertainties only.  The deceleration parameter is
$q_0 = \Omega_M/2 - \Omega_\Lambda$; negative means accelerating.}
\label{tab:results}
\medskip
\begin{tabular}{lccc}
\toprule
Analysis & $\hat\Omega_M$ & $\hat\Omega_\Lambda$ & $\hat q_0$ \\
\midrule
Perlmutter et al.\ (SCP)          & $0.28\pm0.09$ & $0.72\pm0.09$ & $-0.58$ \\
Riess et al.\ (HZT), MLCS         & $0.24\pm0.10$ & $0.76\pm0.10$ & $-0.64$ \\
Riess et al.\ (HZT), template     & $0.20\pm0.10$ & $0.80\pm0.10$ & $-0.70$ \\
\midrule
Planck 2018 (CMB + BAO)           & $0.315\pm0.007$ & $0.685\pm0.007$ & $-0.528$ \\
\bottomrule
\end{tabular}
\end{table}

%======================================================================
\section{What the Supernovae Actually Measured}
\label{sec:fisher}

The Fisher information for \eqref{eq:meas} is
\begin{equation}\label{eq:fisher}
  \mathcal{I}_{jk} = \sum_i \frac{1}{\sigma_i^2}
  \frac{\partial\mu(z_i)}{\partial\theta_j}
  \frac{\partial\mu(z_i)}{\partial\theta_k},
\end{equation}
and inverting it bounds the variance of any unbiased estimator.  The
calculation needs only the redshifts, the error bars, and the model.  It does
not need the magnitudes.  Table~\ref{tab:forecast} reports it for four
designs, each with 18 nearby supernovae and $\mathcal{M}$ profiled out.

\begin{table}[ht]
\centering
\caption{Fisher standard errors for four observing designs, each paired with
18 low-redshift supernovae at $\sigma = 0.17$~mag and profiled over
$\mathcal{M}$.  High-redshift errors are $0.22$~mag, except the last row,
which uses $0.25$.}
\label{tab:forecast}
\medskip
\begin{tabular}{lccccc}
\toprule
High-$z$ design & $n$ & $z$ range & SE($\Omega_M$) & SE($\Omega_\Lambda$)
  & SE($0.8\Omega_M - 0.6\Omega_\Lambda$) \\
\midrule
SCP-like  & 42 & 0.18--0.83 & 0.44 & 0.65 & 0.10 \\
Truncated & 42 & 0.18--0.62 & 0.81 & 0.98 & 0.13 \\
Bunched   & 42 & 0.35--0.55 & 1.30 & 1.50 & 0.12 \\
HZT-like  & 16 & 0.16--0.62 & 1.32 & 1.56 & 0.19 \\
\bottomrule
\end{tabular}
\end{table}

The individual density parameters are not measured.  The SCP-like design gives
a standard error of $0.44$ on $\Omega_M$ and $0.65$ on $\Omega_\Lambda$.  The
first exceeds the estimate itself, $0.28$; the second is $90\%$ of the
estimate $0.72$.  Neither parameter is determined at all in the ordinary
sense.

One combination is measured well.  The forecast standard error on
$0.8\Omega_M - 0.6\Omega_\Lambda$ is $0.10$, and \citet{Perlmutter1999} report
that combination as $-0.2 \pm 0.1$.  A calculation using nothing but redshifts
and error bars reproduces the published error bar.  The observing design fixed
the 1998 constraint; the details of the photometry did little to it.

The reason is a near-degeneracy in \eqref{eq:friedmann}.  Over the observed
range, raising $\Omega_M$ and raising $\Omega_\Lambda$ bend $\mu(z)$ in nearly
opposite ways, with $\partial\mu/\partial\Omega_M \approx -1.1\,
\partial\mu/\partial\Omega_\Lambda$ at $z = 0.5$ and a ratio nearer $-1.5$
when weighted across the whole sample.  Raising $\Omega_M$ by one unit and
$\Omega_\Lambda$ by about $1.5$ therefore leaves the Hubble diagram nearly
unchanged.  \citet{GoobarPerlmutter1995} derived this geometry before the data
existed and designed the survey around it.

\subsection*{Design beats sample size}

The lower rows of Table~\ref{tab:forecast} price the observing strategy.  Keep
the 42 supernovae and truncate the sample at $z = 0.62$ and the standard error
on $\Omega_M$ nearly doubles.  Bunch all 42 near $z = 0.45$ and it triples.
The well-measured combination barely moves in either case, from $0.10$ to
$0.13$ and $0.12$.

Reach in redshift buys separation of the two parameters.  It buys almost
nothing on the combination.  That is why the events beyond $z = 0.6$ were
worth their cost in telescope time: they rotate the confidence ellipse, and
rotation is what tells $\Omega_M$ from $\Omega_\Lambda$.  Where you put the
observations mattered more than how many you took, which is the same lesson
Millikan's spectral range teaches in the photoelectric case.

Here the parallel with Millikan ends, and the difference is the harder half of
the problem.  Millikan's equation had one slope, so removing the contact
voltage left it fully identified.  The Friedmann model has two shape
parameters, and removing $\mathcal{M}$ leaves them tangled with each other.  A
second degeneracy survives the first.  The contact voltage was an instrument
nuisance that better instruments removed; this one is a property of the
geometry, and no supernova measurement removes it.

\subsection*{Where the three sigma came from}

The published significance therefore rests on more than the supernovae.

Measured by the marginal standard error on $\Omega_\Lambda$ alone, the
SCP-like forecast puts $\Omega_\Lambda = 0.72$ just $1.1$ standard errors from
zero.  The physical restriction $\Omega_M \geq 0$ does the rest.  Setting
$\Omega_\Lambda = 0$ and minimising over non-negative $\Omega_M$ raises
$\chi^2$ by $3.8$ in the same forecast, two standard errors.  No non-negative
matter density reproduces the observed curvature with $\Omega_\Lambda = 0$,
and that is what excludes a decelerating universe.  The size of
$\hat\Omega_\Lambda$ on its own does not.

Impose flatness instead and the picture changes.  The same forecast gives a
standard error of $0.070$ on $\Omega_M$, putting Einstein--de Sitter ten
standard errors away.  Flatness is what turns a two-parameter problem into a
one-parameter problem, and it is what the published $\pm 0.09$ describes.

Neither prior is controversial.  The cosmic microwave background supplies
flatness, and its constraint in the $(\Omega_M, \Omega_\Lambda)$ plane runs
nearly orthogonal to the supernovae's, so the two together determine what
neither determines alone.  Both priors are nonetheless external to the
supernova data, and the published error bars do not survive their removal.

Partial identification of this kind will be familiar to economists.  The data
pin down a linear combination of structural parameters; the rest comes from a
restriction imported from outside.  The estimate is only as good as the
restriction, and a careful presentation states the restriction alongside the
number.

%======================================================================
\section{Two Teams, and the Objections}

The two groups shared almost nothing but the answer.  They used different
telescopes and observing strategies.  The High-Z team fitted multi-colour
light-curve templates with a shape parameter; the Berkeley team stretched a
single-colour template along the time axis.  Their $K$-corrections, which
translate a redshifted spectrum back to a common rest-frame filter, came from
different spectral templates.  Their colour corrections used different
estimators.  Their high-redshift samples barely overlapped.

Their best-fit matter densities, $0.28$ and $0.24$, differ by less than half a
standard error.

That agreement tests less than it appears to in one respect and more in
another.  Both numbers assume flatness, and under flatness the supernovae
measure a single parameter, so two teams agreeing on one number with $\pm0.10$
error bars is a weak test.  Against that, a systematic error in any pipeline
step would have had to be reproduced by an independent implementation to
survive.  Replication across methods tests the standardisation.  It does not
test the statistics.

Both teams then had to answer three objections, all of which proposed that
some systematic effect made distant supernovae look faint without any
acceleration.

\emph{Grey dust.}  A screen of intergalactic dust that absorbs all
wavelengths equally would dim distant supernovae without the reddening that
normally betrays dust.  Both teams checked the colour residuals.  Dust of the
required amount would leave a correlation between colour excess and magnitude
offset.  None was found at the level needed.

\emph{Luminosity evolution.}  If supernovae at $z \approx 0.5$ were
intrinsically fainter than nearby ones, because their progenitors were younger
or differently composed, the signal would be spurious.  The teams compared the
decline-rate and colour distributions of near and distant events, and compared
their spectra.  Within measurement error the distant supernovae looked like
nearby ones of the same light-curve shape.

\emph{Selection.}  A magnitude-limited survey preferentially finds
intrinsically bright events at high redshift, which biases the measured
magnitudes the wrong way for acceleration but must still be quantified.  Both
teams simulated their own selection functions and found the effect small.

None of these checks is decisive on its own, and the teams said so.  The
decisive test was the one in the next section, and it took six more years.

%======================================================================
\section{The Test That Came Later}
\label{sec:prediction}

The 1998 fits estimated parameters from the data they fitted.  The test
described here estimated nothing.

Take the flat model with $\Omega_M = 0.28$ and $\Omega_\Lambda = 0.72$, fitted
to supernovae below $z = 0.85$, and ask what it says about $z = 1.5$.
Referred to Einstein--de Sitter it predicts an excess faintness of $+0.41$
magnitudes at $z = 0.5$ and only $+0.74$ at $z = 1.5$.  The curve flattens,
because dark energy dominates late and matter dominates early.

Now take the leading systematic explanation of the same 1998 data.  Grey
intergalactic dust and luminosity evolution both dim distant supernovae by
amounts that grow with distance.  Calibrated to reproduce the observed $+0.41$
at $z = 0.5$, a dimming proportional to $z$ predicts $+1.24$ at $z = 1.5$ and
keeps climbing.

\begin{table}[ht]
\centering
\caption{Excess faintness relative to Einstein--de Sitter under the flat model
fitted in 1998, and under a dimming proportional to $z$ calibrated to match it
at $z = 0.5$.  Per-supernova errors at these redshifts were $0.2$ to $0.3$
magnitudes.}
\label{tab:prediction}
\medskip
\begin{tabular}{lrrr}
\toprule
$z$ & $\Lambda$CDM & Dimming $\propto z$ & Difference \\
\midrule
0.5 & $+0.41$ & $+0.41$ & $0.00$ \\
1.0 & $+0.62$ & $+0.83$ & $-0.20$ \\
1.5 & $+0.74$ & $+1.24$ & $-0.50$ \\
\bottomrule
\end{tabular}
\end{table}

Half a magnitude at $z = 1.5$ is two to three per-supernova standard
deviations, which a handful of events settles.  Referred to an empty universe
instead, the prediction is sharper: $\Lambda$CDM puts supernovae beyond $z
\approx 1.4$ \emph{brighter} than the benchmark, the signature of the
deceleration that preceded the dark-energy era.  No model in which distant
supernovae are dimmed by dust or by evolution produces a change of sign.

\citet{Riess2004} collected the supernovae at $z > 1$ with the Hubble Space
Telescope and found the turnover.  The 1998 parameters had no freedom left to
accommodate the result.  That observation is the out-of-sample test of the
model.  The three sigma of 1998 measured parameters within it.

%======================================================================
\section{What They Could Not Settle}
\label{sec:limits}

Type Ia supernovae measure one function, the luminosity distance $d_L(z)$.
Every quantity above is a functional of that one curve.  Two theories
delivering the same $d_L(z)$ over the observed range are observationally
identical to these data, whatever else they say.

Three classes of theory do so.  A cosmological constant, with equation-of-state
parameter $w = -1$, is one.  A rolling scalar field with $w(z)$ within a few
percent of $-1$ is a second; current data bound $w$ only to about $0.03$,
giving $w = -1.03 \pm 0.03$.  Modified theories of gravity that reproduce the
same expansion history without any dark-energy fluid are a third.  The 1998
data select the shape of $d_L(z)$.  They do not select the physics generating
it.

Distinguishing them needs a different observable.  \citet{Riess2004} and the
later surveys pursued one, the effect of $w$ on $d_L(z)$ at higher redshift.
The growth rate of large-scale structure supplies the other: modified gravity
and dark energy predict different growth rates even when they produce
identical expansion histories, so growth separates them where distances
cannot.  Theory names the discriminating observable.  No one would measure
growth rates to test dark energy without a theory predicting that the two
differ.

And the mechanism remains unexplained.  Quantum field theory predicts a vacuum
energy density some $10^{120}$ times the observed $\Lambda$
\citep{Weinberg1989}.  Riess, Schmidt and Perlmutter measured the curvature of
the Hubble diagram to a few percent.  The discrepancy between that measurement
and the theory of what it measures is the largest in physics.

%======================================================================
\section{The Machine-Learning Reading}

The vocabulary of statistical learning fits this episode closely, and the
preceding note sets out why the translation is more than a conceit
\citep{sargent2025sources}.  Four things a trained network could not deliver
here mark the boundary of what fitting a curve accomplishes.

Start with what the physics supplies.  Friedmann's equations narrow the
admissible curves from all functions $m_B(z)$ to a
three-parameter family: $\mathcal{M}$, $\Omega_M$, $\Omega_\Lambda$.  The
restriction follows deductively from general relativity, the assumption of a
homogeneous and isotropic universe, and the scaling of matter density with
expansion.  Each had been tested independently long before anyone looked for a
high-redshift supernova.

Compare the alternative.  The modest three-layer network priced in the
preceding note carries 593 free parameters against the Friedmann model's
three, and wants some six thousand labelled examples where the Friedmann model
wants four.  Perlmutter had 42 supernovae.  Riess had 50.  The compilations of
the 2010s, with several hundred, fall short by an order of magnitude, and
today's samples of a few thousand by a factor of three.

Give the network enough data and it would eventually trace the curve.  Four
things would still be beyond it.

It could not identify the learned function as $5\log_{10}\tilde d_L(z)$ with
$\tilde d_L$ satisfying a particular integral from general relativity.  The
curve would be a curve.

It could not extract $\Omega_M$ and $\Omega_\Lambda$ from its weights, because
the weights do not decompose that way without the Friedmann equations imposed
as an external constraint.

It could not relate the curvature to the deceleration parameter
$q_0 = \Omega_M/2 - \Omega_\Lambda$, and so could not answer the question
everyone wanted answered.

And it could not predict the $z > 1$ turnover of
Section~\ref{sec:prediction}.  Trained on data
below $z = 1$, it would extrapolate the accelerating trend, because nothing in
its architecture requires a transition.  The observation that settled the 1998
result was a prediction of exactly the kind a function approximator does not
make.

%======================================================================
\section{What This Says About AI in the Sciences}

The preceding note argued that a correctly specified model can turn a problem
needing thousands of observations into one needing five, and that no volume of
data on the wrong observable substitutes for a theory naming the right one.
The supernova episode bears both out.  It also complicates them in a way the
photoelectric case does not, and the complication is the more interesting
half.

Induction did indispensable work here.  Phillips found his relation by fitting
curves to light curves; no theory predicted it, and none derives it now.
Without it the scatter is twice the signal, and a perfectly specified
Friedmann model fitted to unstandardised magnitudes recovers nothing.  Without
Friedmann the standardised magnitudes are a cloud of points with no parameter
in them worth naming.  The discovery required a deductively specified
hypothesis class and an inductively discovered feature, and neither half was
optional.  Anyone arguing that theory beats learning, or that learning will
eventually dispense with theory, has to explain this case away.

The pattern has persisted.  SALT2 and SALT3, on which modern supernova
cosmology runs, are light-curve models trained on data rather than derived
from explosion physics, and their training sets grow with every survey.
Feature extraction in this field has become more machine-learned over time.
The hypothesis class has stayed Friedmann's.  The division of labour is
stable, and it does not fall along the line the usual arguments assume.

On the value of scale the supernova numbers are sharper than the photoelectric
ones, because here the sample size can be held fixed while the design varies.
Every row of Table~\ref{tab:forecast} has the same 42 supernovae.  Bunching
them near $z = 0.45$ triples the standard error on $\Omega_M$ and leaves the
well-measured combination almost untouched.  Only the aim of the telescope
changes across the rows, and the model is what aimed it.

The supernova case then adds a limit with no photoelectric counterpart.
Millikan's slope was identified; the density parameters are not.  Fifty
supernovae fix one combination of the two and leave the other open, so the
published numbers close the gap with $\Omega_M \geq 0$, or with flatness taken
from the microwave background.  Neither restriction came out of the data being
analysed.  A procedure that returns a point estimate and a standard error
while saying nothing about which restriction is carrying them has concealed
the part of the inference a reader most needs to check.  That failure is not
peculiar to neural networks.  Anyone who runs a regression and reports the
output can commit it.

Section~\ref{sec:limits} states the outer limit.  Unlimited supernova data
would not separate a cosmological constant from a rolling scalar field,
because the two imply the same $d_L(z)$.  The separation waits on measurements
of the growth of structure, which exist as a target only because a theory
predicts that the two differ there.

%======================================================================
\section{Summary}

The discovery of cosmic acceleration is a three-parameter weighted nonlinear
regression whose functional form comes from general relativity.  The
regression is straightforward.  Two other things were not.  Standardising the
supernovae by an empirical relation nobody had derived cut the noise by a
factor of $2.7$.  Assembling supernovae at $z = 0.3$ to $0.8$ put observations
where the competing models differ by two or three tenths of a magnitude.
Neither task is statistical.

Under flatness, the estimates of both groups give
\[
  \hat q_0 = \frac{\hat\Omega_M}{2} - \hat\Omega_\Lambda = -0.58,
  \qquad
  \hat\Omega_\Lambda = 0.72 \pm 0.09,
  \qquad
  z_t \approx 0.73 ,
\]
where $z_t$ is the redshift at which deceleration gave way to acceleration.
Without flatness the individual parameters are not determined, and the
inference to acceleration runs through the restriction $\Omega_M \geq 0$.

Fifty supernovae sufficed because the physics had fixed the functional form
and left three parameters.  The fifty could not separate two of those three,
and the published numbers close that gap with a restriction the supernovae do
not supply.  Both facts belong in any account of what the data
established.

\bigskip
\begin{center}\rule{0.4\textwidth}{0.4pt}\end{center}
\smallskip
\noindent\footnotesize
The Fisher forecasts use the observed redshift distributions and quoted
uncertainties, not the individual magnitudes; the parameter estimates are the
two teams' own published values.  The significance figures attributed to
\citet{Riess1998}, the calibration in \eqref{eq:phillips}, and the bound on
$w$ come from secondary sources and should be checked against the originals.

\bigskip
\bibliographystyle{plainnat}
\bibliography{riess_schmidt_perlmutter_ML}

\end{document}
